\documentclass[11pt]{article}

% Packages
\usepackage{kotex}
\usepackage{graphicx}
\usepackage{indentfirst}
\usepackage{setspace}
\usepackage[linguistics]{forest}       % Syntax Tree 패키지 
\usepackage[a4paper, 
            top = 2.5cm, 
            bottom = 2.5cm, 
            left = 2cm, 
            right = 2cm]{geometry}     % 용지 여백 조절 패키지 
\usepackage{fancyhdr}                  % 머리글, 바닥글 조정 패키지
\pagestyle{fancy}
\fancyhf{} % 기존 머리글과 바닥글 초기화
\fancyhead[R]{Syntax(001) - quiz(1)}   % 오른쪽: 문서 제목
\fancyfoot[L]{\nouppercase{\leftmark}} % 왼쪽: 현재 section 제목
\fancyfoot[R]{\thepage}                % 오른쪽: 페이지 번호
\renewcommand{\sectionmark}[1]{%
  \markboth{#1}{}%
}
\renewcommand{\headrulewidth}{0.4pt}   % 머리글 아래쪽 선
\renewcommand{\footrulewidth}{0.4pt}   % 바닥글 위쪽 선
\setlength{\headheight}{14pt}          % 머리글 높이 확보

% 제목 페이지용 바닥글
\fancypagestyle{firstpage}{%
  \fancyhf{}%
  \fancyfoot[L]{1. Tree}%
  \fancyfoot[R]{\thepage}%
  \renewcommand{\headrulewidth}{0pt}%
  \renewcommand{\footrulewidth}{0.4pt}%
}

\setcounter{secnumdepth}{0}            % section 넘버링 제거
\forestset{
  terminal/.style={
    no edge,
    l=3ex,
    for parent={l sep=0pt}
  }
}

% Short form for a lexical category and its terminal.
\forestset{
  lex/.style={append={[#1, terminal]}},
  simple terminals/.style={
    delay={for tree={split option={content}{ }{content,lex}}}
  }
}

% Reusable GPS tree block; default width is 0.48\textwidth.
\NewDocumentEnvironment{gpstree}{O{0.48\textwidth} m +b}
{%
  \begin{minipage}[t]{#1}
  \vspace{0pt}%
  \textbf{#2}\par
  \medskip
  \centering
  \makebox[\linewidth][c]{%
    \begin{forest} simple terminals
      #3
    \end{forest}%
  }
  \end{minipage}%
}
{}

% Metadata
\title{2026 Fall 몽골어 1(001): \\ \medskip \large{laboratory midterm quiz}}
\author{Болорцэцэг  /  1-102}
\date{October 13, 2026}

\section{1. 제시된 자료를 보고, 빈칸을 채우시오.}
\begin{enumerate}
  \item [] 
  \item []
  \item []
\end{enumerate}
\section{2. 제시된 한국어 뜻을 보고 그에 해당하는 몽골어 단어를 적거나, 몽골어 단어를 보고 그것의 한국어 뜻을 적으시오.}
\section{3. 십자말풀이}
\section{4. 자기소개: 몽골어로 자신을 소개할 수 있는 문장을 적으시오. (2문장 이상)}

\end{document}